\chapter*{平面向量}
\subsection*{平面向量的加法与减法}
\begin{enumerate}
    \item 已知向量$\vv{a} = (cos\alpha, sin\alpha), \vv{b} = (cos\beta, sin\beta)$，且$\vv{a} \neq \pm \vv{b}$，那么$\vv{a} + \vv{b}$与$\vv{a} - \vv{b}$的夹角为\blkda{$\dfrac{\pi}{2}$}
    \begin{jieda}
        $(\vv{a} + \vv{b}) \cdot (\vv{a} - \vv{b}) = 0$ \\
        \ps 可以看出$\vv{a}, \vv{b}$是单位向量，构成等腰三角形，而$\vv{a} + \vv{b}$在其角平分线上，故由几何关系可知$(\vv{a} + \vv{b}) \perp (\vv{a} - \vv{b})$
    \end{jieda}
    \item $P$ 是 $\bigtriangleup  {ABC}$ 所在平面上一点，满足 $\left| {\vv{PB} - \vv{PC}}\right|  - \left| {\vv{PB} + \vv{PC} - 2\vv{PA}}\right|  = 0$ ，则 $\bigtriangleup  {ABC}$ 的形状是\dotfill{(\kh{A})}
    \xx{直角三角形}{等腰三角形}{等腰直角三角形}{等边三角形}
    \begin{jieda}
        因为 $\left| {\vv{PB} - \vv{PC}}\right|  - \left| {\vv{PB} + \vv{PC} - 2\vv{PA}}\right|  = 0$ ,
        
        即 $\left| \vv{CB}\right|  - \left| {\left( {\vv{PB} - \vv{PA}}\right)  + \left( {\vv{PC} - \vv{PA}}\right) }\right|  = 0$ ,
        
        即 $\left| \vv{CB}\right|  = \left| {\vv{AB} + \vv{AC}}\right|$ ,
        
        即 $\left| {\vv{AB} - \vv{AC}}\right|  = \left| {\vv{AB} + \vv{AC}}\right|$ ,
        
        所以 ${\left( \vv{AB} - \vv{AC}\right) }^{2} = {\left( \vv{AB} + \vv{AC}\right) }^{2}$ ,
        
        即 ${\vv{AB}}^{2} - 2\vv{AB} \cdot  \vv{AC} + {\vv{AC}}^{2} = {\vv{AB}}^{2} + 2\vv{AB} \cdot  \vv{AC} + {\vv{AC}}^{2}$ ,
        
        所以 $\vv{AB} \cdot  \vv{AC} = 0$ ,所以 $\vv{AB} \bot  \vv{AC}$ ,即 ${AB} \bot  {AC}$ ,
        
        所以 $\bigtriangleup  {ABC}$ 为直角三角形.
    \end{jieda}
\end{enumerate}
\subsection*{定比分点}
\begin{enumerate}
    \item 在$\triangle ABC$中，点$F$为线段$BC$上端点外的任意一点，若$\vv{AF} = x \vv{AB} + 2y \vv{AC}(x > 0, y > 0)$，则$\dfrac{1}{x} + \dfrac{2}{y}$的最小值为\blkda{9}
    \begin{jieda}
        根据等和线结论可知$x + 2y = 1$,因此有$\dfrac{1}{x} + \dfrac{2}{y} = (x+2y)\cdot\left(\dfrac{1}{x} + \dfrac{2}{y}\right) = 1 + 4 + \dfrac{2y}{x} + \dfrac{2x}{y} \geq 5 + 2\sqrt{4} = 9 \left(x = y = \dfrac{1}{3} \text{时取等}\right)$
    \end{jieda}
    \item 已知两点$M(7,8), N(1, -6)$，$P$点是线段$MN$靠近点$M$的三等分点，则点$P$的坐标为\blkda{$\left(5, \dfrac{10}{3}\right)$}
    \begin{jieda}
        $\vv{OP} = \dfrac{2}{3}\vv{OM} + \dfrac{1}{3} \vv{ON} = \left(5, \dfrac{10}{3}\right)$
    \end{jieda}
    \item 已知$A(1,2), B(-4,2)$，$C$满足$\vv{AC} = \lambda \vv{CB}$，$O$为坐标原点，若$|\vv{OC}| = \sqrt{13}$，则$\lambda = $\blkda{4 或$-\dfrac{2}{7}$}.
    \begin{jieda}
        $\vv{OC} = \dfrac{1}{1+\lambda} \vv{OA} + \dfrac{\lambda}{1+\lambda} \vv{OB}$，故$C\left(\dfrac{1-4\lambda}{1+\lambda}, 2\right)$，故$13 = 4 + \left(\dfrac{1-4\lambda}{1+\lambda}\right)^2$，解得$\lambda = 4$
    \end{jieda}
\end{enumerate}
\subsection*{两组三点共线+平面向量基本定理}
\begin{enumerate}
    \item 已知$A,O,B$是不共线的三点，且$C,D$满足$\vv{OC} = 2 \vv{OA}, \vv{OD} = 3 \vv{OB}$，$AD$与$BC$交于点$M$，设$\vv{OM} = x \vv{OA} + y \vv{OB}$，试求$x, y$.
    \begin{jieda}
        $\vv{OM} = x \vv{OA} + y \vv{OB} \Rightarrow \left \{ \begin{array}{l}
                \vv{OM} = \dfrac{x}{2}\vv{OC} + y\vv{OB}  \\
                \vv{OM} = x \vv{OA} + \dfrac{y}{3}\vv{OD} 
        \end{array} \right.$，因为$A, M, D$共线，$C, M, B$共线，根据等和线结论有$\left \{ \begin{array}{l}
                \dfrac{x}{2} + y = 1  \\
                x + \dfrac{y}{3} = 1 
        \end{array} \right.$，解得$x = \dfrac{4}{5}, y = \dfrac{3}{5}$.
    \end{jieda}
    \item 在$\triangle{ABC}$中, 点$P$是$AB$上一点, $\vv{CP}=\dfrac{2}{3}\vv{CA}+\dfrac{1}{3}\vv{CB}$, $Q$是$BC$中点, $AQ$与$CP$交点为$M$, $\vv{CM}=t\vv{CP}$, 则$t=$\blkda{$\dfrac{3}{4}$}
    \begin{jieda}
        $\vv{CM} = \lambda \vv{CA} + (1-\lambda) \vv{CQ} = \lambda \vv{CA} + \dfrac{1-\lambda}{2} \vv{CB}$，因此$\vv{CP} = \dfrac{\lambda}{t} \vv{CA} + \dfrac{1-\lambda}{2t} \vv{CB}$，由平面向量基本定理可推得$\dfrac{1-\lambda}{2t} = \dfrac{1}{2} \cdot \dfrac{\lambda}{t}$，因此$\lambda = \dfrac{1}{2}$，即$\vv{CM} = \dfrac{1}{2}\vv{CA} + \dfrac{1}{4}\vv{CB}$，因此$t = \dfrac{3}{4}$
    \end{jieda}
\end{enumerate}
\subsection*{三点共线系数调整}
\begin{enumerate}
    \item 已知正三角形 ${ABC}$ 的边长为 2，点 $D$ 满足 $\vv{CD} = m\vv{CA} + n\vv{CB}$ ，且 $m > 0,n > 0,{2m} + n = 1$ ，则 $\left| \vv{CD}\right|$ 的取值范围是\blkda[2]{$\left( {1,2}\right)$}.
        \begin{jieda}
            设$E$为$CA$中点，因此$\vv{CD} = 2m\vv{CE} + n\vv{CB}$，因为$m > 0,n > 0,{2m} + n = 1$，所以$D$在线段$BE$上（不含端点），因此$\left| \vv{CD}\right|$ 的取值范围是$\left( {1,2}\right)$.

            \begin{center}
                \includegraphics[width=0.2\textwidth]{bodyxb/src/96_701_1591_168_141_0.jpg}
            \end{center}
        \end{jieda}
    \item (2024$\cdot$徐汇区二模$\cdot$12)如图所示，已知 $\bigtriangleup  {ABC}$ 满足 ${BC} = 8,{AC} =  {3AB},P$ 为 $\bigtriangleup {ABC}$ 所在平面内一点. 定义点集 $D = \left\{  {P\left| {\vv{AP} = {3\lambda }\vv{AB} + }\right. }\right.  \left. {\dfrac{1 - \lambda }{3}\vv{AC},\lambda  \in  \mathbf{R}}\right\}$. 若存在点 ${P}_{0} \in  D$ ,使得对任意 $P \in  D$ ,满足 $\left| \vv{AP}\right|  \geq  \left| \vv{A{P}_{0}}\right|$ 恒成立，则 $\left| \vv{A{P}_{0}}\right|$ 的最大值为\blkda[2]{$3$}.

    \begin{center}
    \includegraphics[width=0.2\textwidth]{bodyxb/src/135_1120_548_330_144_0.jpg}
    \end{center}
    \begin{jieda}
        取$M$满足$\vv{AM} = 3\vv{AB}$，$N$满足$\vv{AN} = \dfrac{1}{3}\vv{AC}$，故集合$D$为直线$MN$，$\left| \vv{AP}\right|  \geq  \left| \vv{A{P}_{0}}\right|$ 恒成立即$AP_0 \bot MN$，根据对称性可知$A$到$BC$的距离和$A$到$MN$距离相等，又因为$A$在半径为$3$的阿氏圆上，因此$\left| \vv{A{P}_{0}}\right|$ 的最大值为$3$.
        \begin{center}
        \includegraphics[width=0.2\textwidth]{bodyxb/src/99_578_509_242_178_0.jpg}
        \end{center}
    \end{jieda}
\end{enumerate}
\subsection*{平面向量平行}
\begin{enumerate}
    \item 已知$\vv{a}, \vv{b}$是非零不共线向量，向量$8\vv{a} - k \vv{b}$与$-k\vv{a} + \vv{b}$共线，则$k = $ \blkda{$\pm 2\sqrt{2}$}
        \begin{jieda}
            根据向量共线充要条件可以得到$8 = -\lambda k, -k = \lambda$，即$k^2 = 8$，解得$k = \pm 2\sqrt{2}$
        \end{jieda}
    \item (2025$\cdot$宝山区二模$\cdot$13)已知向量 $\vv{a} = \left( {1,x}\right)$ ， $\vv{b} = \left( {3,1}\right)$ ，若 $\vv{a}//\vv{b}$ ，则 $x$ 的值为 \dotfill{(\kh{D})}
    \xx{$-3$}{$3$}{$- \dfrac{1}{3}$}{$\dfrac{1}{3}$}

    \item 已知向量$\vv{a} = (1, 2), \vv{b} = (2, -2), \vv{c} = (1, \lambda)$，若$\vv{c} // (2\vv{a} + \vv{b})$，则$\lambda = $\blkda{$\dfrac{1}{2}$}
    \begin{jieda}
        $(2\vv{a} + \vv{b}) = (4, 2)$，因为$\vv{c} // (2\vv{a} + \vv{b})$，所以根据平面向量平行的充要条件的坐标表示有$4\lambda = 2$，故$\lambda = \dfrac{1}{2}$
    \end{jieda}
\end{enumerate}
\subsection*{数量积的定义}
\begin{enumerate}
    \item 设$\vv{a}, \vv{b}$是任意的向量，则下列命题正确的是\blkda{\ding{175}}:
    \begin{dingautolist}{172}
        \item 若$\vv{a} \neq \vv{0}, \vv{a} \cdot \vv{b} = 0$，则$\vv{b} = \vv{0}$
        \item 若$\vv{a} \cdot \vv{b} = \vv{b} \cdot \vv{c}$，则$\vv{a} = \vv{c}$
        \item $(\vv{a} \cdot \vv{b}) \vv{c} = \vv{a} (\vv{b} \cdot \vv{c})$
        \item $\vv{a} \cdot [(\vv{a} \cdot \vv{c})\cdot \vv{b} - (\vv{a} \cdot \vv{b})\cdot \vv{c}] = 0$
    \end{dingautolist}

    \item 在$\triangle{ABC}$中, 若$BC=5$, $AC=4$, $\angle{C}=\ang{45}$, 则$\vv{BC}\cdot\vv{CA}=$\blkda{$-10\sqrt{2}$}. 
    \item ${AB}$ 为圆 $O$ 的一条弦，且 $\left| {AB}\right|  = 2$ ，则 $\vv{AB} \cdot  \vv{AO}$ 的值为\blkda{2}
    \begin{jieda}
        取弦 ${AB}$ 的中点 $M$ ,连接 ${OM}$ ,根据圆的垂径定理,可得 ${OM} \bot  {AB}$ ,如图.
        
        \begin{center}
        \includegraphics[width=0.15\textwidth]{bodyxb/src/0_242_740_257_258_0.jpg}
        \end{center}
        
        
        因为 $\left| {AB}\right|  = 2$ ,所以 $\left| {AM}\right|  = 1$ .
        
        根据向量数量积的几何意义:
        
        $\vv{AB} \cdot  \vv{AO} = \vv{AO} \cdot  \vv{AB} = \left| {AM}\right| \left| {AB}\right|  = 1 \times  2 = 2$
    \end{jieda}
    \item (2024$\cdot$徐汇区一模$\cdot$9)在 $\bigtriangleup  {ABC}$ 中， ${AC} = {BC},{P}_{1}$ 、 ${P}_{2}$ 、 ${P}_{3}$ 为边 ${AB}$ 上的点，且 $8\left| \vv{{P}_{1}B}\right|  = 4\left| \vv{{P}_{2}B}\right|  =  2\left| \vv{{P}_{3}B}\right|  = \left| \vv{AB}\right|  = 8$ ,设 ${I}_{k} = \vv{{P}_{k}B} \cdot  \vv{{P}_{k}C}\left( {k = 1,2,3}\right)$ ,则 ${I}_{1} - {I}_{2} + {I}_{3} =$ \blkda[2]{1}.
    \begin{jieda}
        如图所示,由题意得 ${P}_{3}B = 4$ , ${P}_{2}B = 2,{P}_{1}B = 1$ ,且 ${P}_{3}C \bot  {P}_{3}B$ . 从而 ${I}_{1} = \vv{{P}_{1}B} \cdot  \vv{{P}_{1}C} =  - {P}_{1}B \cdot  {P}_{1}{P}_{3} =  - 3;  {I}_{2} = \vv{{P}_{2}B} \cdot  \vv{{P}_{2}C} =  - {P}_{2}B \cdot  {P}_{2}{P}_{3} =  - 4;  {I}_{3} = \vv{{P}_{3}B} \cdot  \vv{{P}_{3}C} = 0$ ; 因此 ${I}_{1} - {I}_{2} + {I}_{3} = 1$.

        \begin{center}
        \includegraphics[width=0.2\textwidth]{bodyxb/src/97_608_776_212_175_0.jpg}
        \end{center}
    \end{jieda}
    \item (2024$\cdot$崇明区二模$\cdot$11)已知 $A$ 、 $B$ 、 $C$ 是半径为 1 的圆上的三个不同的点，且 $\left| \vv{AB}\right|  = \sqrt{3}$ ，则 $\vv{AB} \cdot  \vv{AC}$ 的最小值是\blkda[2]{$\dfrac{3}{2} - \sqrt{3}$}.
    \begin{jieda}
        $\vv{AC}$在$\vv{AB}$上数量投影最小即可.
    \end{jieda}
    \item 若平面内单位向量$\vv{a_1}, \vv{a_2}, ..., \vv{a_n}$满足对于任意的$1 \leq i < j \leq n$，都有$\vv{a_i} \cdot \vv{a_j} < \dfrac{1}{2}$，则正整数$n$的最大值为\blkda{5}
    \begin{jieda}
        对于单位向量，数量积小于$\dfrac{1}{2}$即夹角大于$\dfrac{\pi}{3}$，由于任意两个单位向量之间的夹角都要满足大于$\dfrac{\pi}{3}$，因此最多可以有$5$个单位向量.
    \end{jieda}
\end{enumerate}
\subsection*{利用数量积求夹角、求模}
\begin{enumerate}
    \item (2024$\cdot$青浦区二模$\cdot$2)已知向量 $\vv{a} = \left( {-1,1}\right) ,\vv{b} = \left( {3,4}\right)$ ，则 $\langle \vv{a},\vv{b}\rangle  =$ \blkda[2]{$\arccos \dfrac{\sqrt{2}}{10}$}.
    \item (2024$\cdot$静安区二模$\cdot$4)若单位向量 $\vv{a}$ 、 $\vv{b}$ 满足 $\vv{a}\bot \vv{b}$ ，则 $\left| {\vv{a} - \sqrt{3}\vv{b}}\right|  =$ \blkda[2]{2}.
\end{enumerate}
\subsection*{平面向量垂直}
\begin{enumerate}
    \item (2025$\cdot$杨浦区一模$\cdot$6)已知 $\left| \vv{a}\right|  = \left| \vv{b}\right|  = 1$ ，若$(2\vv{a} - \vv{b})\bot \vv{b}$ ，则向量 $\vv{a}$ 与 $\vv{b}$ 的夹角的余弦值为\blkda[2]{$\dfrac{1}{2}$}.
    \begin{jieda}
        设向量 $\vv{a}$ 与 $\vv{b}$ 的夹角为 $\theta \left( {\theta  \in  \left\lbrack  {0,\pi }\right\rbrack  }\right)$ ,若 $\left( {2\vv{a} - \vv{b}}\right)  \bot  \vv{b}$ ,则 $\left( {2\vv{a} - \vv{b}}\right)  \cdot  \vv{b} = 0$ ,所以 $2\vv{a} \cdot  \vv{b} - {\vv{b}}^{2} = 2\left| \vv{a}\right| \left| \vv{b}\right| \cos \theta  - {\left| \vv{b}\right| }^{2}  = 2\cos \theta  - 1 = 0$ ,可得 $\cos \theta  = \dfrac{1}{2}$.
    \end{jieda}
    \item 两个非零向量$\vv{a},\vv{b}$互相垂直的充要条件是\dotfill(\kh{C})
    \xx
    {$\vv{a}\cdot\vv{b}=\abs{\vv{a}}\cdot\abs{\vv{b}}$}
    {$\vv{a}\cdot\vv{b}=\vv{0}$}
    {$\abs{\vv{a}+\vv{b}}=\abs{\vv{a}-\vv{b}}$}
    {$(\vv{a}+\vv{b})\cdot(\vv{a}-\vv{b})=0$}
\end{enumerate}
\subsection*{数量投影和投影向量}
\begin{enumerate}
    \item (2024$\cdot$奉贤区二模$\cdot$8)已知向量 $\vv{a} = \left( {1,1}\right) ,\vv{b} = \left( {2, - 1}\right)$ ，则 $\vv{b}$ 在 $\vv{a}$ 方向上的投影向量为\blkda[2]{$\left( {\dfrac{1}{2},\dfrac{1}{2}}\right)$}.
    \item (2025$\cdot$青浦区一模$\cdot$12)已知 $A$ 、 $B$ 、 $C$ 是单位圆上任意不同三点，则 $\dfrac{\vv{AB} \cdot  \vv{AC}}{\left| \vv{AB}\right| }$ 的取值范围是\blkda[2]{$\left( {-1,2}\right)$}.
    \begin{jieda}
        $\dfrac{\vv{AB} \cdot  \vv{AC}}{\left| \vv{AB}\right| }$即$\vv{AC}$在$\vv{AB}$上的数量投影，以$\vv{AB}$为$x$轴正方向，设$A(x_0, y_0), x_0 \in (-1, 0)$
        
        固定$x_0$，则$\vv{AC}$在$x$轴上的数量投影取值范围是$(-1-x_0, 1-x_0)$，因此$\dfrac{\vv{AB} \cdot  \vv{AC}}{\left| \vv{AB}\right| } \in \left( {-1,2}\right)$

        \twotu{\includegraphics[width=0.4\textwidth]{bodyxb/src/98_1022_1632_180_169_0.jpg}}{\includegraphics[width=0.4\textwidth]{bodyxb/src/98_1275_1632_189_169_0.jpg}}
    \end{jieda}
    \item (2024$\cdot$普陀区二模$\cdot$9)若向量 $\vv{a}$ 在向量 $\vv{b}$ 上的投影为 $\dfrac{1}{3}\vv{b}$ ，且 $\left| {3\vv{a} - \vv{b}}\right|  = \left| {\vv{a} + \vv{b}}\right|$ ，则 $\cos \langle \vv{a},\vv{b}\rangle  =$ \blkda[2]{$\dfrac{\sqrt{3}}{3}$}.
    \begin{jieda}
        因为 $\vv{a}$ 在 $\vv{b}$ 上的投影为 $\dfrac{1}{3}\vv{b}$ ,所以 $\dfrac{\vv{a} \cdot  \vv{b}}{\left| \vv{b}\right| } \cdot  \dfrac{\vv{b}}{\left| \vv{b}\right| } = \dfrac{1}{3}\vv{b}$ , 则 $\dfrac{\vv{a} \cdot  \vv{b}}{{\left| \vv{b}\right| }^{2}} = \dfrac{1}{3}$ ,即 $\left| \vv{b}\right|  = \sqrt{3\vv{a} \cdot  \vv{b}}$ ,又 $\left| {3\vv{a} - \vv{b}}\right|  = \left| {\vv{a} + \vv{b}}\right|$ ,平方得 $8{\vv{a}}^{2} = 8\vv{a} \cdot  \vv{b}$ ,则 $\left| \vv{a}\right|  = \sqrt{\vv{a} \cdot  \vv{b}}$ ,即 $\cos \langle \vv{a},\vv{b}\rangle  = \dfrac{\vv{a} \cdot  \vv{b}}{\left| \vv{a}\right| \left| \vv{b}\right| } =  \dfrac{\vv{a} \cdot  \vv{b}}{\sqrt{3\vv{a} \cdot  \vv{b}}\sqrt{\vv{a} \cdot  \vv{b}}} = \dfrac{\sqrt{3}}{3}$.
    \end{jieda}
\end{enumerate}
\subsection*{基的思想}
\begin{enumerate}
    \item 在$\triangle ABC$中，$\angle A = \ang{90}$，$AB = 1, AC = 2$，设$P, Q$满足$\vv{AP} = \lambda \vv{AB}, \vv{AQ} = (1-\lambda)\vv{AC}, \lambda \in \mathbf{R}$，若$\vv{BQ} \cdot \vv{CP} = -2$，则$\lambda = $ \blkda{$\dfrac{2}{3}$}
    \begin{jieda}
        选择$\vv{AB}, \vv{AC}$为一组基底 \\
        $\vv{BQ} \cdot \vv{CP} = (\vv{BA} + \vv{AP}) \cdot (\vv{CA} + \vv{AP}) = ((1-\lambda)\vv{AC} - \vv{AB}) \cdot (\lambda \vv{AB} - \vv{AC}) = -\lambda|\vv{AB}|^2 - (1-\lambda)|\vv{AC}|^2 = 3\lambda - 4 = -2$，故$\lambda = \dfrac{2}{3}$
    \end{jieda}

    \item 在$\triangle ABC$中，$O$为$BC$中点，点$M,N$分别在边$AB, AC$上，且$AM = 6, MB = 4, AN = 4, NC = 3$，$\vv{OM} \cdot \vv{ON} = 0$， 则$cos A = $\blkda{$\dfrac{3}{8}$}
    \begin{jieda}
        首先考虑以$\vv{AB}, \vv{AC}$为基底表示$\vv{OM}, \vv{ON}$：\\
        $\left \{ \begin{array}{l}
                \vv{OM} = \vv{OA} + \vv{AM} = -\dfrac{1}{2}(\vv{AB} + \vv{AC}) + \dfrac{3}{5} \vv{AB}  \\
                \vv{ON} = \vv{OA} + \vv{AN} = -\dfrac{1}{2}(\vv{AB} + \vv{AC}) + \dfrac{4}{7} \vv{AC}
        \end{array} \right .$ ，因此$(\dfrac{1}{10}\vv{AB} - \dfrac{1}{2}\vv{AC}) \cdot (-\dfrac{1}{2}\vv{AB} + \dfrac{1}{14}\vv{AC}) = 0$ \\
        即$\dfrac{9}{35} \vv{AB} \cdot \vv{AC} - \dfrac{1}{20}|\vv{AB}|^2 - \dfrac{1}{28}|\vv{AC}|^2 = 0$，代入$|\vv{AB}| = 10, |\vv{AC}| = 7$，得到$cos A = \dfrac{3}{8}$
    \end{jieda}

\end{enumerate}

\subsection*{数量积的范围（最值）问题与主元归一}
\begin{enumerate}

    \item 在$\triangle ABC$中，$\angle BAC = \dfrac{\pi}{3}$，$\vv{AD} = 2\vv{DB}$，$P$为$CD$上一点，且满足$\vv{AP} = m\vv{AC} + \dfrac{1}{3}\vv{AB}$，若$\vv{AB} \cdot \vv{AC} = 4$，则$|\vv{AP}|$的最小值为\blkda{2}
    \begin{jieda}
        $\vv{AD} = 2\vv{DB} \Rightarrow \vv{AD} = \dfrac{2}{3}\vv{AB}$，故$\vv{AP} = m \vv{AC} + \dfrac{1}{2}\vv{AD}$，根据等和线结论可知$m = \dfrac{1}{2}$ \\
        因此$|\vv{AP}|^2 = (\dfrac{1}{2}\vv{AC} + \dfrac{1}{3}\vv{AB})^2 = \dfrac{1}{4}|\vv{AC}|^2 + \dfrac{4}{3} + \dfrac{1}{9}|\vv{AB}|^2$\\
        根据$\vv{AB} \cdot \vv{AC} = 4$，$\angle BAC = \dfrac{\pi}{3}$知$|\vv{AB}| \cdot |\vv{AC}| = 8$ \\
        设$|\vv{AC}| = t$，则$|\vv{AP}|^2 = \dfrac{t^2}{4} + \dfrac{4}{3} + \dfrac{64}{9t^2} \geq 4$（当$t = \dfrac{4\sqrt{3}}{3}$时取等），故$|\vv{AP}| = 2$
    \end{jieda}
    \item (2026$\cdot$上海春季卷$\cdot$10)在 $\bigtriangleup  {ABC}$ 中， $D$ 、 $E$ 在边 ${BC}$ 上，且 $\vv{BD} = \vv{DE} = \vv{EC}$ ， $\left| \vv{AD}\right|  = 1$ ， $\vv{AD}$ 与 $\vv{AE}$ 所成的夹角为 $\dfrac{\pi }{3}$ ，则 $\vv{AB} \cdot  \vv{AC}$ 的最大值为\blkda{$- \dfrac{39}{32}$}.
    
    \begin{center}
    \includegraphics[width=0.2\textwidth]{bodyxb/src/4_240_432_342_185_0.jpg}
    \end{center}
    \begin{jieda}
        $\because \vv{AB} = \vv{AD} + \vv{DB} = \vv{AD} + \vv{ED} = \vv{AD} + \vv{EA} + \vv{AD} = 2\vv{AD} - \vv{AE}$
        
        $\vv{AC} = \vv{AD} + \vv{DC} = \vv{AD} + 2\vv{DE} = \vv{AD} + 2\left( {\vv{DA} + \vv{AE}}\right)  =  - \vv{AD} + 2\vv{AE}$
        
        $\therefore \vv{AB} \cdot  \vv{AC} = \left( {2\vv{AD} - \vv{AE}}\right)  \cdot  \left( {-\vv{AD} + 2\vv{AE}}\right)  =  - 2{\left| \vv{AD}\right| }^{2} + 5\vv{AD} \cdot  \vv{AE} - 2{\left| \vv{AE}\right| }^{2}$
        
        $\because \left| \vv{AD}\right|  = 1,\vv{AD}$ 与 $\vv{AE}$ 所成的夹角为 $\dfrac{\pi }{3}$
        
        $\therefore \vv{AB} \cdot  \vv{AC} =  - 2 + \dfrac{5}{2} \cdot  \left| \vv{AE}\right|  - 2{\left| \vv{AE}\right| }^{2}$
        
        令 $\left| \vv{AE}\right|  = t > 0$ ,则 $\vv{AB} \cdot  \vv{AC} =  - 2{t}^{2} + \dfrac{5}{2}t - 2 =  - 2{\left( t - \dfrac{5}{8}\right) }^{2} - \dfrac{39}{32},t > 0$
        
        当 $\left| \vv{AE}\right|  = \dfrac{5}{8}$ 时, $\vv{AB} \cdot  \vv{AC}$ 的最大值为 $- \dfrac{39}{32}$ .
    \end{jieda}

    
    \item 在$\param ABCD$中，$E$为$CD$中点，$AB = 2$，若$AE \perp BD$，求$\vv{AC} \cdot \vv{AB}$的取值范围.
    \begin{jieda}
        以$\vv{AC}$和$\vv{AB}$为基底入手，表示$\vv{AE}$和$\vv{BD}$并使其数量积为0建立等式. \\
        $\left \{\begin{array}{l}
                \vv{AE} = \vv{AC} + \vv{CE} =  \vv{AC} - \dfrac{1}{2}\vv{AB}  \\
                \vv{BD} = \vv{BA} + \vv{AC} + \vv{CD} = \vv{AC} - 2\vv{AB} 
        \end{array} \right. $ \\
        因此有$\vv{AE} \cdot \vv{BD} = (\vv{AC} - \dfrac{1}{2}\vv{AB}) \cdot (\vv{AC} - 2\vv{AB}) = |\vv{AC}|^2 - \dfrac{5}{2} \vv{AC} \cdot \vv{AB} + |\vv{AB}|^2 = 0$\\
        设$|\vv{AC}| = x, \vv{AC}, \vv{AB}$的夹角为$\theta$，则$x^2 - 5x \cdot cos\theta + 4 = 0$，即$x + \dfrac{4}{x} = 5 cos\theta$，根据余弦函数的值域和对勾函数的性质可知$x \in (1,4)$时方程有解. 故$\vv{AC} \cdot \vv{AB} = 2x \cdot cos\theta = \dfrac{2(4+x^2)}{5} \in (2, 8)$
    \end{jieda}
    
    \item 在正方形$ABCD$中，边长为2，$M$在$BC$边上，$N$在$CD$边上，满足$MN=\sqrt{2}$，求$\vv{AM}\cdot \vv{AN}$的取值范围\blkda{$[4, 8-2\sqrt{2}]$}
    \begin{jieda}
        \begin{dingautolist}{172}
            \item 方法1
            
            以A为原点，$\vv{AB}$为$x$轴正方向建系，设$\angle MNC = \theta, \theta \in \left[0, \dfrac{\pi}{2}\right]$，则$M(2, 2-\sqrt{2}sin\theta), N(2-\sqrt{2}cos\theta, 2)$，故$\vv{AM}\cdot \vv{AN} = 4-2\sqrt{2}cos\theta + 4 - 2\sqrt{2}sin\theta = 8 - 4sin\left(\theta + \dfrac{\pi}{4}\right) \in [4, 8-2\sqrt{2}]$
            \item 方法2
            
            取$MN$中点$E$，则$\vv{AM} \cdot \vv{AN} = \left|\vv{AE}\right|^2 - \left(\dfrac{\sqrt{2}}{2}\right)^2$

            $\left|\vv{CE}\right| = \dfrac{\sqrt{2}}{2}$，故$E$在以$C$为圆心的圆弧上，因此$\left|\vv{AE}\right|^2 \in \left[\dfrac{9}{2}, \dfrac{17}{2}-2\sqrt{2}\right]$

            故$\vv{AM}\cdot \vv{AN} \in [4, 8-2\sqrt{2}]$
        \end{dingautolist}
    \end{jieda}
\end{enumerate}
\subsection*{极化恒等式}
\begin{enumerate}
    \item 正六边形边长为2，点$O$为正六边形的中心，若点$M$在正六边形的外接圆上运动，点$A, B$是半径为1的小圆$O$的直径，则$\vv{MA} \cdot \vv{MB} = $\blkda{3}
        \begin{jieda}
            $\vv{MA} \cdot \vv{MB} = (\vv{MO} + \vv{OA}) \cdot (\vv{MO} + \vv{OB}) = |\vv{MO}|^2 + \vv{MO} \cdot (\vv{OA}+\vv{OB}) + \vv{OA} \cdot \vv{OB} = 3$ \\
            \ps 该结论（$O$为$\vv{AB}$中点，则$\vv{MA} \cdot \vv{MB} = |\vv{MO}|^2 - |\vv{OA}|^2$）称为极化恒等式.
        \end{jieda}
        \item 如图，在$\triangle ABC$中，$D$是$BC$的中点，$E, F$是$A, D$上的两个三等分点，$\vv{BA} \cdot \vv{CA} = 4$，$\vv{BF} \cdot \vv{CF} = -1$，则$\vv{BE} \cdot \vv{CE}$的值是 \blkda{$\dfrac{7}{8}$} 
        \begin{center}
            \includegraphics[width=0.2\linewidth]{bodyxb/src/pingmianxiangliang_34q.png}
        \end{center}
        \begin{jieda}
            选择$\vv{BD}, \vv{DA}$为基，这样计算数量积的时候可以因为$\vv{BD} = -\vv{CD}$实现消项的目的.\\
            $\vv{BA} \cdot \vv{CA} = (\vv{BD} + \vv{DA}) \cdot (\vv{CD} + \vv{DA}) = -|\vv{BD}|^2 + |\vv{DA}|^2 = 4$ \\
            $\vv{BF} \cdot \vv{CF} = (\vv{BD} + \vv{DF}) \cdot (\vv{CD} + \vv{DF}) = -|\vv{BD}|^2 + |\vv{DF}|^2 = -|\vv{BD}|^2 + \dfrac{1}{9}|\vv{DA}|^2 = -1$ \\
            故$|\vv{DA}|^2 = \dfrac{45}{8}$ \\
            $\vv{BE} \cdot \vv{CE} = (\vv{BD} + \vv{DE}) \cdot (\vv{CD} + \vv{DE}) = -|\vv{BD}|^2 + |\vv{DE}|^2 = -|\vv{BD}|^2 + \dfrac{4}{9}|\vv{DA}|^2 = 4 - \dfrac{5}{9} \cdot \dfrac{45}{8} = \dfrac{7}{8}$
            \ps 本题利用极化恒等式也能解决
        \end{jieda} 

        \item 如图所示,正方形 ${ABCD}$ 边长为 6,圆 $D$ 的半径为 $1,E$ 是圆 $D$ 上任意一点,则 $\vv{AE} \cdot  \vv{CE}$ 的最小值为\blkda{$1 - 6\sqrt{2}$}.
        
        \begin{center}
        \includegraphics[width=0.2\textwidth]{bodyxb/src/1_240_198_307_304_0.jpg}
        \end{center}
        \begin{jieda}
            取$AC$中点$M$，$\vv{AE}\cdot\vv{CE} = |\vv{EM}|^2 - |\vv{CM}|^2$，最小值为 $1 - 6\sqrt{2}$ .
        \end{jieda}
\end{enumerate}
\subsection*{几何性质确定-建系}
\begin{enumerate}
    \item 在边长为1的正三角形$ABC$中, 若$\vv{BC}=\vv{a}$, $\vv{CA}=\vv{b}$, $\vv{AB}=\vv{c}$, 则$\vv{a}\cdot\vv{b}+\vv{b}\cdot\vv{c}+\vv{c}\cdot\vv{a}=$\blkda{$-\dfrac{3}{2}$}
    \item 已知向量$\vv{a}, \vv{b}, \vv{c}$满足$|\vv{a}| = |\vv{b}| = 1, |\vv{c}| = \sqrt{2}$，且$\vv{a} + \vv{b} + \vv{c} = \vv{0}$，则$cos<\vv{a} - \vv{c}, \vv{b} - \vv{c}> = $ \blkda{$\dfrac{4}{5}$}
    \item (2024$\cdot$闵行区一模$\cdot$10)若平面上的三个单位向量 $\vv{a}$ 、 $\vv{b}$ 、 $\vv{c}$ 满足 $\left| {\vv{a} \cdot  \vv{b}}\right|  = \dfrac{1}{2}$ ， $\left| {\vv{a} \cdot  \vv{c}}\right|  = \dfrac{\sqrt{3}}{2}$ ，则 $\vv{b} \cdot  \vv{c}$ 的所有可能的值组成的集合为\blkda[2]{$\left\{  {0,\dfrac{\sqrt{3}}{2}, - \dfrac{\sqrt{3}}{2}}\right\}$}.
    \begin{jieda}
        不妨设$\vv{a} = (1, 0)$，则$\vv{b} = \left(\pm \dfrac{1}{2}, \pm \dfrac{\sqrt{3}}{2}\right), \vv{c} = \left(\pm \dfrac{\sqrt{3}}{2}, \pm \dfrac{1}{2}\right)$，故$\vv{b} \cdot  \vv{c} \in \left\{  {0,\dfrac{\sqrt{3}}{2}, - \dfrac{\sqrt{3}}{2}}\right\}$
    \end{jieda}
    \item (2025$\cdot$上海秋季卷$\cdot$12)已知 $f\left( x\right)  = \left\{  \begin{array}{ll} 1,&x > 0, \\  0,&x = 0, \\   - 1,&x < 0. \end{array}\right. \vv{a},\vv{b},\vv{c}$ 是平面内三个不同的单位向量. 若 $f\left( {\vv{a} \cdot  \vv{b}}\right)  + f\left( {\vv{b} \cdot  \vv{c}}\right)  + f\left( {\vv{c} \cdot  \vv{a}}\right)  = 0$ ，则 $\left| {\vv{a} + \vv{b} + \vv{c}}\right|$ 的取值范围是\blkda[2]{$\left( {1,\sqrt{5}}\right)$}.
    \begin{jieda}
        若 $f\left( {\vv{a} \cdot  \vv{b}}\right)  = f\left( {\vv{b} \cdot  \vv{c}}\right)  = f\left( {\vv{c} \cdot  \vv{a}}\right)  = 0$ ,则 $\vv{a} \cdot  \vv{b} = \vv{b} \cdot  \vv{c} = \vv{c} \cdot  \vv{a} = 0$ ,
        
        又三个向量均为平面内的单位向量,故向量 $\vv{a},\vv{b},\vv{c}$ 两两垂直,显然不成立; 故 $\left\{  {f\left( {\vv{a} \cdot  \vv{b}}\right) ,f\left( {\vv{b} \cdot  \vv{c}}\right) ,f\left( {\vv{c} \cdot  \vv{a}}\right) }\right\}   = \{  - 1,0,1\}$ .
        
        不妨设 $\left\{  \begin{array}{l} f\left( {\vv{a} \cdot  \vv{b}}\right)  = 1 \\  f\left( {\vv{b} \cdot  \vv{c}}\right)  = 0 \\  f\left( {\vv{c} \cdot  \vv{a}}\right)  =  - 1 \end{array}\right.$ ,则 $\vv{a} \cdot  \vv{b} > 0,\vv{b} \cdot  \vv{c} = 0,\vv{c} \cdot  \vv{a} < 0$ ,
        
        不妨设 $\vv{b} = \left( {1,0}\right) ,\vv{c} = \left( {0,1}\right) ,\vv{a} = \left( {\cos \theta ,\sin \theta }\right) ,\theta  \in  \lbrack 0,{2\pi })$ ,
        
        则 $\left\{  \begin{array}{l} \vv{a} \cdot  \vv{b} = \cos \theta  > 0 \\  \vv{c} \cdot  \vv{a} = \sin \theta  < 0 \end{array}\right.$ ,则 $\theta  \in  \left( {\dfrac{3}{2}\pi ,{2\pi }}\right)$ ,
        
        则 $\left| {\vv{a} + \vv{b} + \vv{c}}\right|  = \left| \left( {1 + \cos \theta ,1 + \sin \theta }\right) \right|  = \sqrt{{\left( 1 + \cos \theta \right) }^{2} + {\left( 1 + \sin \theta \right) }^{2}} = \sqrt{3 + 2\cos \theta  + 2\sin \theta }  = \sqrt{3 + 2\sqrt{2}\sin \left( {\theta  + \dfrac{\pi }{4}}\right) }$ ,
        
        由 $\theta  \in  \left( {\dfrac{3}{2}\pi ,{2\pi }}\right) ,\theta  + \dfrac{\pi }{4} \in  \left( {\dfrac{7}{4}\pi ,\dfrac{9}{4}\pi }\right)$ ,
        
        则 $\sin \left( {\theta  + \dfrac{\pi }{4}}\right)  \in  \left( {-\dfrac{\sqrt{2}}{2},\dfrac{\sqrt{2}}{2}}\right) ,2\sqrt{2}\sin \left( {\theta  + \dfrac{\pi }{4}}\right)  \in  \left( {-2,2}\right)$
        
        故 $\left| {\vv{a} + \vv{b} + \vv{c}}\right|  \in  \left( {1,\sqrt{5}}\right)$ .
    \end{jieda}
\end{enumerate}

\subsection*{圆上动点}
\begin{enumerate}
    \item (2025$\cdot$杨浦区二模$\cdot$15)已知 $A$ 、 $B$ 、 $C$ 是单位圆上的三个点，若 $\left| \vv{AB}\right|  = \sqrt{2}$ ，则 $\vv{AB} \cdot \vv{BC}$ 的最大值为 \dotfill{(\kh{D})}
    \xx{$\sqrt{2}$}{$1 + \dfrac{\sqrt{2}}{2}$}{$\sqrt{2} + 1$}{$\sqrt{2} - 1$}

    \item 已知等边三角形$ABC$的边长为$2\sqrt{3}$，点$P$是其外接圆上的动点，则$\vv{PA} \cdot \vv{PB} + \vv{PA} \cdot \vv{PC}$的取值范围为\blkda{$[0, 8]$}
    \begin{jieda}
        设$\triangle ABC$外接圆圆心为$O$，以$\vv{OA}$的方向为$x$轴正方向，设$A(2, 0)$, $B(-1, \sqrt{3})$, $C(-1, -\sqrt{3})$，建立平面直角坐标系。设$P(2cos\theta, 2sin\theta)$，由此\\
        $\vv{PA} = (2-2cos\theta, -2sin\theta), \vv{PB} = (-1-2cos\theta, \sqrt{3}-2sin\theta), \\
        \vv{PC} = (-1-2cos\theta, -\sqrt{3}-2sin\theta), \vv{PB} + \vv{PC} = (-2-4cos\theta, -4sin\theta)$，\\
        故$\vv{PA} \cdot \vv{PB} + \vv{PA} \cdot \vv{PC} = \vv{PA} \cdot (\vv{PB}+\vv{PC}) = (2-2cos\theta) \cdot (-2-4cos\theta) + (-2sin\theta) \cdot (-4sin\theta) = -4 - 4cos\theta + 8cos^2\theta + 8sin^2\theta = 4 - 4cos\theta \in [0, 8]$
    \end{jieda}
\end{enumerate}

\subsection*{模的几何意义-向量差的模}
\begin{enumerate}
    \item 已知平面向量 $\vv{a}$ 、 $\vv{b}$ 满足 $\left| \vv{a}\right|  = 3,\left| \vv{b}\right|  = 4,\vv{a} \cdot  \vv{b} = 4$ ，若平面向量 $\vv{c}$ 满足 $\left| {\vv{c} - \vv{b}}\right|  = 1$ ，则 $\left| {\vv{c} - \vv{a}}\right|$ 的最大值为\blkda[2]{$\sqrt{17} + 1$}.
    \begin{jieda}
        以$\vv{b}$为$x$轴正方向建系，$A(1, 2\sqrt{2}), B(4, 0)$，$|AB| = \sqrt{17}$，$C$在以$B$为圆心，$1$为半径的圆上，$|AC| \leq \sqrt{17} + 1$
    \end{jieda}
    \item (2024$\cdot$崇明区一模$\cdot$11)已知不平行的两个向量 $\vv{a}$ 、 $\vv{b}$ 满足 $\left| \vv{a}\right|  = 1$ ， $\vv{a} \cdot  \vv{b} = \sqrt{3}$. 若对任意的 $t \in  \mathbf{R}$ ， 都有 $\left| {\vv{b} - t\vv{a}}\right|  \geq  2$ 成立，则 $\left| \vv{b}\right|$ 的最小值等于\blkda[2]{$\sqrt{7}$}.
    \begin{jieda}
        向量 $\vv{a},\vv{b}$ 满足 $\left| \vv{a}\right|  = 1$，$\vv{a} \cdot  \vv{b} = \sqrt{3} \Longrightarrow  \left| \vv{a}\right| \left| \vv{b}\right| \cos \theta  = \sqrt{3} \Longrightarrow  \left| \vv{b}\right| \cos \theta  = \sqrt{3}$ ,

        对任意的 $t \in  \mathbf{R}$ ,都有 $\left| {\vv{b} - t\vv{a}}\right|  \geq  2 \Longrightarrow  \left| {\vv{b} - t\vv{a}}\right|  = \left| \vv{b}\right| \sin \theta  \geq  2$ 成立,当 ${BA} \bot  {OA}$ 时 $m = \left| {\vv{b} - t\vv{a}}\right|$ 取得最小值,如图所示, $\left\{  {\begin{array}{l} \left| \vv{b}\right| \cos \theta  = \sqrt{3} \\  \left| \vv{b}\right| \sin \theta  \geq  2 \end{array} \Longrightarrow  \left| +\right| }\right.  {\left| \vv{b}\right| }^{2} \geq  7,\left| \vv{b}\right|  \geq  \sqrt{7}.$

        \begin{center}
        \includegraphics[width=0.2\textwidth]{bodyxb/src/100_1341_2020_236_140_0.jpg}
        \end{center}
    \end{jieda}
    \item 已知平面内的向量$\vv{a}, \vv{b}, \vv{c}, \vv{e}$满足$|\vv{a}| = 4, |\vv{e}| = 1, \vv{a} \perp \vv{e}, |\vv{b} - \vv{a}| = 1$，且有$|\vv{c} - \lambda \vv{e}| \geq |\vv{c} - \vv{e}|$对任意$\lambda \in \mathbf{R}$恒成立，则$|\vv{c} - \vv{b}|$的最小值为\blkda{0}.
    \begin{jieda}
        以$\vv{a}$为$x$轴正方向，$\vv{e}$为$y$轴正方向建立平面直角坐标系，设$\vv{a} = (4, 0), \vv{e} = (0, 1)$，$\vv{b} = (4+cos\theta, sin\theta)$，其中$\theta$为$\vv{b} - \vv{a}$与$x$轴正方向夹角，$|\vv{c} - \lambda \vv{e}| \geq |\vv{c} - \vv{e}|$即$\vv{c} = (m, 1)$，故$|\vv{c} - \vv{b}| \geq 0$（$\vv{c} = \vv{b} = (4, 1)$时取得最小值）
    \end{jieda}
    \item (2025$\cdot$上海春季卷$\cdot$12)在平面中， ${\vv{e}}_{1}$ 和 ${\vv{e}}_{2}$ 是互相垂直的单位向量，向量 $\vv{a}$ 满足 $\left|\vv{a} - 4\vv{e_1}\right|  = 2$ ，向量 $\vv{b}$ 满足 $\left|\vv{b} - 6\vv{e_2}\right|  = 1$ ，则 $\vv{b}$ 在 $\vv{a}$ 方向上的数量投影的最大值\blkda[2]{4}.
    \begin{jieda}
        设 $\vv{a} = \vv{OA},\vv{b} = \vv{OB}$ ,在坐标系中,不妨设 $4{\vv{e}}_{1} = \left( {0,4}\right)  = \vv{O{O}_{1}},6{\vv{e}}_{2} = \left( {6,0}\right)  = \vv{O{O}_{2}}$ ,则 $\left|\vv{OA} -  \vv{O{O}_{1}}\right|  = \left| \vv{{O}_{1}A}\right| = 2, \left|\vv{OB} - \vv{O{O}_{2}}\right| = \left|\vv{{O}_{2}B} \right| = 1$
        
        如图,显然 ${OA}$ 与圆 ${O}_{1}$ 相切在右侧时, $\vv{b}$ 在 $\vv{a}$ 方向上的数量投影最大, 当 ${BH} \bot  {OA}$ ,且 ${BH}$ 与圆 ${O}_{2}$ 相切时, $\vv{b}$ 在 $\vv{a}$ 方向上的数量投影 $\left| {OH}\right|$ 最大,由题意,得 $\angle {O}_{1}{OA} = {30}^{ \circ  } = \angle {HTO} \Longrightarrow  \left| {{O}_{2}T}\right|  = 2$ ,因为 $\bigtriangleup {HOT} \backsim  \bigtriangleup B{O}_{2}T$ ,所以 $\left| {OH}\right|  = 4$，故 $\vv{b}$ 在 $\vv{a}$ 方向上的数量投影的最大值为4.
        \begin{center}
            \includegraphics[width=0.4\textwidth]{bodyxb/src/95_516_1093_304_208_0.jpg}
        \end{center}
    \end{jieda}
\end{enumerate}

\subsection*{平面向量的外积与三角形面积}
\begin{enumerate}
    \item 已知向量 $\vv{AB} = \left( {1,2}\right)$ ， $\vv{AC} = \left( {4, - 2}\right)$ ，则 $\bigtriangleup  {ABC}$ 的面积为\blkda{5}.
    \item (2020$\cdot$崇明区二模$\cdot$11)在$\triangle ABC$中，$\vv{AB} = (\sqrt{3}cosx, cosx), \vv{AC} = (cosx, sinx)$，则$\triangle ABC$面积最大值为\blkda{$\dfrac{3}{4}$}
    \begin{jieda}
        \newline
        $S_{\triangle ABC} = \dfrac{|\sqrt{3}sinx \cdot cosx - cos^2x|}{2} = \dfrac{1}{2}\left| \dfrac{\sqrt{3}}{2}sin2x - \dfrac{1+cos2x}{2}\right| = \dfrac{1}{2} \left | sin(2x-\dfrac{\pi}{6}) - \dfrac{1}{2} \right| \leq \dfrac{3}{4}$\\
        （当$x = k\pi - \dfrac{\pi}{6}, k \in \mathbf{Z}$时取最大值.）
    \end{jieda}
\end{enumerate}

\subsection*{平面向量的内积与三角形面积}
\begin{enumerate}
    \item 在$\triangle ABC$中，$BC = 2\sqrt{6}, S_{\triangle ABC} = \dfrac{\sqrt{2}}{2}(\vv{AB} \cdot \vv{AC})$，则$\triangle ABC$的外接圆半径为\blkda{$3$}
    \begin{jieda}
        首先根据三角形面积建立等式：$\dfrac{\sqrt{2}}{2}bc cosA = \dfrac{1}{2}bc sinA$，因此$tanA = \sqrt{2}$，故$cosA = \dfrac{\sqrt{3}}{3}, sinA = \dfrac{\sqrt{6}}{3}$。根据正弦定理可知$2R = \dfrac{a}{sinA} = \dfrac{2\sqrt{6}}{\dfrac{\sqrt{6}}{3}} = 6$，故$R = 3$
    \end{jieda}

    \item 在 $\bigtriangleup  {ABC}$ 中， $\vv{AB} = \left( {\cos \theta ,\sin \theta }\right)$ ， $\vv{BC} = \left( {3\sin \theta ,3\cos \theta }\right)$ . 若 $\vv{AB} \cdot  \vv{BC} = 2$ ，则 $\bigtriangleup  {ABC}$ 的面积为\blkda{$\dfrac{\sqrt{5}}{2}$}.
    \begin{jieda}
        $\left| \vv{AB}\right|  = \sqrt{{\cos }^{2}\theta  + {\sin }^{2}\theta } = 1,\left| \vv{BC}\right|  = \sqrt{9{\sin }^{2}\theta  + 9{\cos }^{2}\theta } = 3$ ,
        
        则 $\vv{AB} \cdot  \vv{BC} =  - \vv{BA} \cdot  \vv{BC} =  - \left| \vv{BA}\right|  \cdot  \left| \vv{BC}\right| \cos B =  - 3\cos B = 2$ ,所以 $\cos B =  - \dfrac{2}{3}$ ,
        
        所以 $\sin B = \sqrt{1 - {\cos }^{2}B} = \dfrac{\sqrt{5}}{3}$ .
        
        所以 ${S}_{\bigtriangleup {ABC}} = \dfrac{1}{2}\left| \vv{AB}\right|  \cdot  \left| \vv{BC}\right|  \cdot  \sin B = \dfrac{1}{2} \times  1 \times  3 \times  \dfrac{\sqrt{5}}{3} = \dfrac{\sqrt{5}}{2}$ .
    \end{jieda}
    \item 在 $\bigtriangleup  {ABC}$ 中, $\vv{AB} = \left( {\sqrt{3}\cos x,\cos x}\right) ,\vv{AC} = \left( {\cos x,\sin x}\right)$ . 对任意的实数 $t$ ,恒有 $\left| {\vv{AB} - t\vv{AC}}\right|  \geq  \left| \vv{BC}\right|$ ,则 $\bigtriangleup {ABC}$ 面积的最大值为\blkda{$\dfrac{\sqrt{3}}{2}$}.
    \begin{jieda}
        $\because \left| {\vv{AB} - t\vv{AC}}\right|  \geq  \left| \vv{BC}\right|$ ，故点 $B$ 到直线 ${AC}$ 的最短距离为 ${BC}$ ， $\therefore {BC}\bot {AC}$ . $\because \left| \vv{AB}\right|  = \sqrt{4{\cos }^{2}x} \leq  2,\left| \vv{AC}\right|  = 1$ ， $\therefore$ 由勾股定理可得 ${BC} = \sqrt{A{B}^2 - A{C}^2} \leq  \sqrt{4 - 1} = \sqrt{3}$ ， 故 $\bigtriangleup {ABC}$ 面积 $S = \dfrac{1}{2}{BC} \cdot  {AC} \leq  \dfrac{\sqrt{3}}{2}$.
    \end{jieda}
\end{enumerate}

\subsection*{平面向量与三角形的四心}
\begin{enumerate}
    \item 在$\triangle ABC$中，已知$G$为$\triangle ABC$的重心，内角$A,B,C$的对边长度分别为$a, b, c$，且$56a \vv{GA} + 40b \vv{GB} + 35c\vv{GC} = \vv{0}$,则$B = $\blkda{$\dfrac{\pi}{3}$}
    \begin{jieda}
        因为$G$为$\triangle ABC$的重心，因此$56a = 40 b = 35c$，即$a = \dfrac{5}{8}c, b = \dfrac{7}{8}c$，根据余弦定理有$cosB = \dfrac{a^2 + c^2 - b^2}{2ac} = \dfrac{\dfrac{25}{64} + 1 - \dfrac{49}{64}}{\dfrac{5}{4}} = \dfrac{1}{2}$，故$B = \dfrac{\pi}{3}$
    \end{jieda}
    \item 平面内不重合的点$A_1, A_2, A_3, B_1, B_2, B_3, B_4$，若$|\vv{A_1B_i} + \vv{A_2B_i} + \vv{A_3B_i}| = i$，其中$i = 1, 2, 3, 4$，则$|\vv{B_1B_2}| + |\vv{B_2B_3}| + |\vv{B_3B_4}|$的最大值与最小值之和为\blkda{6}
    \begin{jieda}
        设$G$为$\triangle A_1A_2A_3$的重心，则$|\vv{A_1B_i} + \vv{A_2B_i} + \vv{A_3B_i}| = 3|\vv{GB_i}|$，因此$|\vv{B_1B_2}| + |\vv{B_2B_3}| + |\vv{B_3B_4}|$取最小值的时候四个点和$G$共线，顺次排列，最小值为$\dfrac{1}{3}[1+1+1] = 1$；$|\vv{B_1B_2}| + |\vv{B_2B_3}| + |\vv{B_3B_4}|$取最大值的时候四个点和$G$共线，间隔排列，最大值为$\dfrac{1}{3}[(1+2)+(2+3)+(3+4)] = 5$，故最大值和最小值之和为6.
    \end{jieda}

    \item 在锐角$\triangle{ABC}$中，$cosB = \dfrac{1}{2}$，$O$为$\triangle{ABC}$的外心
    \begin{enumerate}
        \item 若$\vv{BO} = x\vv{BA} + y \vv{BC}$，求$x+y$的最大值
        \item 若$b = \sqrt{3}$，求$|3\vv{OB} + \vv{OA} + 2\vv{OC}|$的取值范围
    \end{enumerate}
    \begin{jieda}
        \begin{enumerate} 
            \item 
            \begin{enumerate}
                \item $\vv{BO} = x\vv{BA} + y \vv{BC}$，故$\left \{ \begin{array}{l}
                    \vv{BO} \cdot \vv{BA} = x|\vv{BA}|^2 + y \vv{BC}\cdot \vv{BA} \\
                    \vv{BO} \cdot \vv{BC} = x\vv{BA} \cdot \vv{BC} + y |\vv{BC}|^2
                \end{array}
                \right.$，\\
                即$\left \{ \begin{array}{l}
                    \dfrac{|\vv{BA}|^2}{2} = x|\vv{BA}|^2 + y \vv{BC}\cdot \vv{BA} \\
                    \dfrac{|\vv{BC}|^2}{2} = x\vv{BA} \cdot \vv{BC} + y |\vv{BC}|^2
                \end{array}
                \right.$，设$|\vv{BA}| = c, |\vv{BC}| = a$，则有\\
                $\left \{ \begin{array}{l}
                    \dfrac{c^2}{2} = x \cdot c^2 + \dfrac{y}{2} \cdot ac \\
                    \dfrac{a^2}{2} = \dfrac{x}{2} \cdot ac + y \cdot a^2
                \end{array}
                \right.$，即$\left \{ \begin{array}{l}
                    \dfrac{c}{2} = x \cdot c + \dfrac{y}{2} \cdot a \\
                    \dfrac{a}{2} = \dfrac{x}{2} \cdot c + y \cdot a
                \end{array}
                \right.$，解得$\left \{ \begin{array}{l}
                    x = \dfrac{2c-a}{3c} \\
                    y = \dfrac{2a-c}{3a}
                \end{array}
                \right.$，故$x+y \leq \dfrac{2}{3}$，此时$x = y = \dfrac{1}{3}$
                \item $\vv{BO} = x(\vv{BO} + \vv{OA}) + y(\vv{BO} + \vv{OC})$，故$(x+y-1)\vv{OB} = x\vv{OA} + y \vv{OC}$，两侧平方有$(x+y-1)^2|\vv{OB}|^2 = x^2|\vv{OA}|^2+y^2|\vv{OC}|^2+2xy|\vv{OB}|\cdot|\vv{OC}|cos\angle BOC$，即$(x+y-1)^2 = x^2+y^2-xy$，即$(x+y-1)^2 -(x+y)^2 = -3xy$，即$2(x+y) = 3xy+1$，根据基本不等式有$2(x+y) \leq \dfrac{3}{4}(x+y)^2+1$（$x=y$时取等号），故$x+y\leq\dfrac{2}{3}$或$x+y\geq2$（舍）因此$x+y$的最大值为$\dfrac{2}{3}$，此时$x = y = \dfrac{1}{3}$
            \end{enumerate}
            \ps 根据等和线的结论可以进行验证
            \item 
            \begin{enumerate}
                \item $cosB = \dfrac{1}{2}$，故$sinB = \dfrac{\sqrt{3}}{2}$，根据正弦定理可知外接圆半径$R = 1$，$|3\vv{OB} + \vv{OA} + 2\vv{OC}|^2 = 9|\vv{OB}|^2+|\vv{OA}|^2+4|\vv{OC}|^2 + 6\vv{OA} \cdot \vv{OB} + 4\vv{OA} \cdot \vv{OC} + 12\vv{OB} \cdot \vv{OC} = 14 + (6cos\angle AOB -2 + 12cos\angle BOC)$，设$\angle BOC$为$\theta, \theta \in \left(\dfrac{\pi}{3}, \pi\right)$，则原式$ = 12+12cos\theta + 6 cos\left(\dfrac{4\pi}{3} - \theta\right) = 12 + 6\sqrt{3}sin\left(\dfrac{\pi}{3} - \theta\right) \in [12-6\sqrt{3}, 12)$ \\
                因此$|3\vv{OB} + \vv{OA} + 2\vv{OC}| \in [3-\sqrt{3}, 2\sqrt{3})$
                \item $cosB = \dfrac{1}{2}$，故$sinB = \dfrac{\sqrt{3}}{2}$，根据正弦定理可知外接圆半径$R = 1$，故以$O$为原点，$\vv{AC}$方向为$x$轴正方向建系，有$A\left(-\dfrac{\sqrt{3}}{2}, -\dfrac{1}{2}\right), C\left(\dfrac{\sqrt{3}}{2}, -\dfrac{1}{2}\right)$因为$\triangle{ABC}$为锐角三角形，因此$B(cos\theta, sin\theta), \theta \in \left(\dfrac{\pi}{6}, \dfrac{5\pi}{6}\right)$，故$3\vv{OB} + \vv{OA} + 2\vv{OC} = \left(3cos\theta+\dfrac{\sqrt{3}}{2}, 3sin\theta-\dfrac{3}{2}\right)$，故$|3\vv{OB} + \vv{OA} + 2\vv{OC}| = \sqrt{12+3\sqrt{3}cos\theta-9sin\theta} = \sqrt{12+6\sqrt{3}\left(\dfrac{1}{2}cos\theta - \dfrac{\sqrt{3}}{2}sin\theta\right)} = \sqrt{12+6\sqrt{3}sin\left(\dfrac{\pi}{6}-\theta\right)} \in [3-\sqrt{3}, 2\sqrt{3})$
                \end{enumerate}
        \end{enumerate}
        \ps 分析对比
        \begin{dingautolist}{172}
            \item 第一问方法1是需要植入的模式，遇到三角形顶点为公共起点的一组基表示外心的式子，$\vv{BO} = x\vv{BA} + y \vv{BC}$，应考虑同时点乘，利用中垂线性质以构建关于系数的二元一次方程，目标是求解两个系数; \\
            而方法2是利用了外心到三个顶点距离相等，但需注意平方前后对于求解系数并非充要，平方操作扩大了解的范围，类似于$A = \sqrt{B}$平方后会扩大$A$的范围. \\
            \\
            从图像角度理解：$2(x+y) = 3xy+1$可化为$y = \dfrac{2x-1}{3x-2}$，为双曲线形式，而原本$(x, y)$只对应双曲线左下方的一支，但平方操作扩大了其范围，对应了两支. \\
            从图像角度线性规划理解求最值的过程，设动直线$x+y = k$，取最值对应了双曲线相切的情形. \\ 
            \\
            从等和线的角度来理解问题：像$\vv{BO} = x\vv{BA} + y \vv{BC}$这样共起点的一组基表示的向量，当其系数和为定值的时候，终点的运动轨迹是一条平行于这一组基的直线. \\
            设$\vv{BO}$终点$O$所在的等和线为$l_1$，$AC$所在直线为$l_2$. 设$B$到$l_2$距离为$h$，易知当$h$大于$l_1, l_2$之间距离$d$时，系数和为正数，否则为负数. 系数和为正时，$B$到两直线距离之比即$O$所在等和线的系数和，即$x+y = \dfrac{h-d}{h}$，因此$h$越大，$x+y$越大，因此当$\triangle ABC$为等腰三角形时取得$x+y$的最大值.
            \item 第二问利用平方操作则是合适的，因为此时对于求模长的最值平方前后是充要的，但无论是采用平方的方式还是建系的方式求解，关键都在于不能忽略锐角三角形的条件，确定好对应的变量的取值范围.
        \end{dingautolist}
    \end{jieda}


    \item 已知在$\triangle{ABC}$中，点$O$是$\triangle{ABC}$的外心，若$|\vv{AB}| = 4$，$\vv{AO} \cdot \vv{BC} = 4$，$\vv{AB} \cdot \vv{AC} = 12\sqrt{2}$，则$\triangle{ABC}$的面积为\blkda{$2\sqrt{6}$}
    \begin{jieda}
        \begin{enumerate}
            \item 方法1 \\
            设$BC$边中点为$D$，因此$\vv{AD} \cdot \vv{BC} = \vv{AO} \cdot \vv{BC} = 4$，进一步以$\vv{AB}, \vv{AC}$为基进行拆解，得到$\dfrac{1}{2}(\vv{AB}+\vv{AC}) \cdot (\vv{AC} - \vv{AB}) = 4$
            \item 方法2 \\
            $\vv{AO} \cdot \vv{BC} = \vv{AO} \cdot (\vv{AC} - \vv{AB}) = \dfrac{|\vv{AC}|^2}{2} - \dfrac{|\vv{AB}|^2}{2} = 4$
        \end{enumerate}
        故$|\vv{AC}|^2 - |\vv{AB}|^2 = 8$，因此$|\vv{AC}| = 2\sqrt{6}$\\
        故$cosA = \dfrac{\vv{AB} \cdot \vv{AC}}{|\vv{AB}| \cdot |\vv{AC}|} = \dfrac{\sqrt{3}}{2}$，因此$sinA = \dfrac{1}{2}$，故$S = \dfrac{1}{2} \times 4 \times 2\sqrt{6} \times \dfrac{1}{2} = 2\sqrt{6}$
    \end{jieda}

    \item 已知$A(0, 1)$和$B(-3, 4)$，若$C$在$\angle AOB$的角平分线上，且$|\vv{OC}| = 2$，则$\vv{OC}$的坐标为\blkda[3]{$\left(-\dfrac{\sqrt{10}}{5}, \dfrac{3\sqrt{10}}{5}\right)$}
    \begin{jieda}
        设$D(0, 5)$，满足$|\vv{OD}| = |\vv{OB}|$，因为$C$在$\angle AOB$的角平分线上，即在$\angle DOB$的角平分线上所以存在实数$\lambda$满足$\vv{OC} = \lambda \cdot (\vv{OB} + \vv{OD})$，因此可设$C(-k, 3k), k>0$，又因为$|\vv{OC}| = 2$，所以$k = -\dfrac{\sqrt{10}}{5}$，因此$\vv{OC} = \left(-\dfrac{\sqrt{10}}{5}, \dfrac{3\sqrt{10}}{5}\right)$
    \end{jieda}

    \item 若$|\vv{AB}| = 4$，$\vv{AC} = 3\vv{CB}$，平面内一点$P$满足$\dfrac{\vv{PA} \cdot \vv{PC}}{|\vv{PA}|} = \dfrac{\vv{PB} \cdot \vv{PC}}{|\vv{PB}|}$，则$sin\angle PAB$的最大值为\blkda{$\dfrac{1}{3}$}
    \begin{jieda}
        $\dfrac{\vv{PA} \cdot \vv{PC}}{|\vv{PA}|} = \dfrac{\vv{PB} \cdot \vv{PC}}{|\vv{PB}|}$即$cos\angle CPA = cos\angle CPB$，即$PC$为$\angle APB$的角平分线. \\
        根据角平分线定理，以及$\vv{AC} = 3\vv{CB}$可知$\dfrac{|\vv{PA}|}{|\vv{PB}|} = 3$ \\
        根据大角对大边可知$\angle PAB$为锐角（或零角，此时不构成三角形），因此其正弦值的最大值对应余弦值最小值，设$x = |\vv{PB}|$根据余弦定理有$cos\angle PAB = \dfrac{9x^2 + 16 - x^2}{24x} = \dfrac{1}{3} \cdot \left(x + \dfrac{2}{x}\right) \geq \dfrac{2\sqrt{2}}{3}$（当$x = \sqrt{2}$时取等号），因此$sin\angle PAB$的大值为$\dfrac{1}{3}$
    \end{jieda}

    \item 在$\triangle ABC$中，$\vv{AB} \perp \vv{AC}$，点$P$在线段$BC$上，且$\vv{AP} \cdot \vv{AC} = \vv{AP} \cdot \vv{AB} = 4$，则$|\vv{AB} + \vv{AC} + \vv{AP}|$的最小值为\blkda{6}
    \begin{jieda}
        根据$\vv{AP} \cdot \vv{AC} = \vv{AP} \cdot \vv{AB} = 4$可知$\vv{AP} \cdot (\vv{AC} - \vv{AB}) = 0$，即$\vv{AP} \perp \vv{BC}$，因此$\vv{AP} \cdot \vv{AB} = |\vv{AP}|^2 = 4$，即$|\vv{AP}| = 2$ \\
        以$P$为原点，$\vv{PA}$为$x$轴正方向建立平面直角坐标系，故$A(2, 0)$，设$B(0, b), C(0, c)$，由$\vv{AB} \perp \vv{AC}$有$bc = -4$ \\
        $\vv{AB} + \vv{AC} + \vv{AP} = (-6, b+c)$，$|\vv{AB} + \vv{AC} + \vv{AP}| = \sqrt{36 + (b+c)^2} \geq \sqrt{36 + b^2 + \dfrac{16}{b^2} - 8} \geq 6$（当$b = 2$时取等）
    \end{jieda}
\end{enumerate}

\subsection*{费马点}
\begin{enumerate}
    \item 我们知道, 三角形中存在诸多特殊位置的点, 并且这些特殊点都具备一定的特殊性质.意大利学者托里拆利在研究时发现: 在三角形的三边分别向其外侧作等边三角形, 这三个等边三角形的外接圆交于一点 $T$ , 该点 $T$ 即称为托里拆利点 (以下简称“ $T$ 点”). 通过研究发现三角形中的“ $T$ 点”满足到三角形三个顶点的距离和 ${TA} + {TB} + {TC}$ 最小. 当 $\bigtriangleup  {ABC}$ 的三个内角均小于120 \textdegree 时,使得 $\angle {AOB} = \angle {BOC} = \angle {COA} = {120}^{ \circ  }$ 的点 $O$ 即为“ $T$ 点”，当 $\vartriangle  {ABC}$ 有一个内角大于或等于120 \textdegree 时，最大内角的顶点为“ $T$ 点”. 试用以上知识解决下面问题: 已知 $\bigtriangleup  {ABC}$ 的内角 $A,B,C$ 所对的边分别为 $a,b,c$ .
    \begin{enumerate}
        \item 若 $\sqrt{3}b - c\sin A = \sqrt{3}a\cos C$ ,则\ding{172}求 $A$，\ding{173}若 ${bc} = 4$ ，设点 $P$ 为 $\bigtriangleup  {ABC}$ 的“ $T$ 点”，求 $\vv{PA} \cdot  \vv{PB} + \vv{PB} \cdot  \vv{PC} + \vv{PC} \cdot  \vv{PA}$
        \item 若 $a\cos B - b\cos A = c$ ，设 $P$ 点为 $\bigtriangleup  {ABC}$ 的“ $T$ 点”， $\left| \vv{PB}\right|  + \left| \vv{PC}\right|  = {2t}\left| \vv{PA}\right|$ ，求实数 $t$ 的最小值.
    \end{enumerate}
    \begin{jieda}
        \begin{enumerate}
            \item \begin{dingautolist}{172}
                \item 在 $\bigtriangleup  {ABC}$ 中，由正弦定理得 $\sqrt{3}\sin B - \sin C\sin A = \sqrt{3}\sin A\cos C$ ，

                    $\because B = \pi  - \left( {A + C}\right)$ ,有 $\sin B = \sin \left( {A + C}\right)$ ,

                    $\sqrt{3}\left( {\sin A\cos C + \cos A\sin C}\right)  - \sin C\sin A = \sqrt{3}\sin A\cos C$ ,

                    $\sqrt{3}\cos A\sin C = \sin C\sin A,\therefore \sin C \neq  0,\therefore \tan A = \sqrt{3}$ ,又 $A \in  \left( {0,\pi }\right) ,\therefore A = \dfrac{\pi }{3}$
                \item 由 \ding{172} 知 $A = \dfrac{\pi }{3}$ ，则 $\bigtriangleup  {ABC}$ 的三个角都小于 ${120}^{ \circ  }$ ，

                    由 “ $T$ 点” 定义知: $\angle {APB} = \angle {BPC} = \angle {APC} = {120}^{ \circ  }$ ,

                    设 $\left| \vv{PA}\right|  = x,\left| \vv{PB}\right|  = y,\left| \vv{PC}\right|  = z$ ,由 ${S}_{\bigtriangleup {APB}} + {S}_{\bigtriangleup {BPC}} + {S}_{\bigtriangleup {APC}} = {S}_{\bigtriangleup {ABC}}$ 得

                    $\dfrac{1}{2}{xy} \cdot  \dfrac{\sqrt{3}}{2} + \dfrac{1}{2}{yz} \cdot  \dfrac{\sqrt{3}}{2} + \dfrac{1}{2}{xz} \cdot  \dfrac{\sqrt{3}}{2} = \dfrac{1}{2} \times  4 \times  \dfrac{\sqrt{3}}{2}$ ,整理得 ${xy} + {yz} + {xz} = 4$ ,

                    所以 $\vv{PA} \cdot  \vv{PB} + \vv{PB} \cdot  \vv{PC} + \vv{PA} \cdot  \vv{PC} = {xy} \cdot  \left( {-\dfrac{1}{2}}\right)  + {yz} \cdot  \left( {-\dfrac{1}{2}}\right)  + {xz} \cdot  \left( {-\dfrac{1}{2}}\right)  =  - \dfrac{1}{2} \times  4 =  - 2$.

            \end{dingautolist}
            \item 由 $a\cos B - b\cos A = c$ ，结合正弦定理 $\sin A\cos B - \sin B\cos A = \sin C$

            有 $\sin \left( {A - B}\right)  = \sin C,\because A,B,C$ 均为三角形内角, $\therefore A - B = C$ 或 $A - B + C = \pi$ (舍),即 $A = B + C = \pi  - A,\therefore A = \dfrac{\pi }{2}$ , 由点 $P$ 为 $\bigtriangleup  {ABC}$ 的“ $T$ 点”，得 $\angle {APB} = \angle {BPC} = \angle {APC} = {120}^{ \circ  }$

            设 $\left| {PB}\right|  = m\left| {PA}\right| ,\left| {PC}\right|  = n\left| {PA}\right| ,\left| {PA}\right|  = x,\left( {m > 0,n > 0,x > 0}\right)$ ,

            ${\left| AB\right| }^{2} = {x}^{2} + {m}^{2}{x}^{2} - {2m}{x}^{2}\cos \dfrac{2\pi }{3} = \left( {{m}^{2} + m + 1}\right) {x}^{2},$

            ${\left| AC\right| }^{2} = {x}^{2} + {n}^{2}{x}^{2} - {2n}{x}^{2}\cos \dfrac{2\pi }{3} = \left( {{n}^{2} + n + 1}\right) {x}^{2},$

            ${\left| BC\right| }^{2} = {m}^{2}{x}^{2} + {n}^{2}{x}^{2} - {2mn}{x}^{2}\cos \dfrac{2\pi }{3} = \left( {{m}^{2} + {n}^{2} + {mn}}\right) {x}^{2}$ ,

            相加得 ${\left| AC\right| }^{2} + {\left| AB\right| }^{2} = {\left| BC\right| }^{2}$ ,得 $\left( {{n}^{2} + n + 1}\right) {x}^{2} + \left( {{m}^{2} + m + 1}\right) {x}^{2} = \left( {{m}^{2} + {n}^{2} + {mn}}\right) {x}^{2}$，整理得 $m + n + 2 = {mn}$ 
            
            由 $\left| {PB}\right|  + \left| {PC}\right|  = {2t}\left| {PA}\right|$ ,得 $m + n = {2t}$

            于是 $m + n + 2 = {mn} \leq  {\left( \dfrac{m + n}{2}\right) }^{2}$（当且仅当 $m = n$ ,即 $m = n = 1 + \sqrt{3}$ 时取等号），即${t}^{2} - {2t} - 2 \geq  0$ ,而 $t > 0$ ,解得 $t \geq  1 + \sqrt{3}$ ,所以实数 $t$ 的最小值为 $1 + \sqrt{3}$ .
        \end{enumerate}
    \end{jieda}
\end{enumerate}